problem:  
Let \((M^n,F,dV_F)\) be a smooth, forward-complete **Finsler manifold** equipped with the **Busemann–Hausdorff measure**, and assume its **weighted Ricci curvature** satisfies the Ohta–Sturm lower bound  
\[
\mathrm{Ric}_N^F \ge (N-1)K, \qquad K \in \mathbb{R}, \; N>n.
\]
Let \(\Sigma=\partial \Omega \subset M\) be a smooth, closed, mean-convex hypersurface bounding a compact domain \(\Omega\), and denote by \(H_F\) the **Finsler mean curvature** associated to the anisotropic second fundamental form (the divergence of the Chern-normalized unit normal).  

For standard Riemannian spaces, the **Heintze–Karcher inequality** controls the integral \(\int_{\Sigma}\frac{1}{H}\,d\mu\) by the enclosed volume, with equality characterizing constant curvature space forms.  
In Minkowski and Euclidean spaces, analogous *anisotropic* Heintze–Karcher-type inequalities hold, but only for \(K=0\).  

**Problem (Curvature–sharp Finsler–Heintze–Karcher conjecture):**  
Does there exist a sharp comparison inequality of the form
\[
\int_{\Sigma} \frac{1}{H_F} \, d\mu_F \;\geq\; \Phi_{K,N,F}\big(\mathrm{Vol}_F(\Omega)\big),
\]
where
- \(\Phi_{K,N,F}\) is an explicit monotone model function determined by the \(N\)-dimensional Finsler space form of constant flag curvature \(K\),  
- equality holds if and only if \((M,F)\) is (locally) isometric to the corresponding Finsler space form (i.e., of constant flag curvature \(K\) and constant S-curvature), and  
- the inequality extends to possibly non-reversible Finsler spaces satisfying the \(CD(K,N)\) condition?

Equivalently, can one interpret such a curvature–sharp anisotropic Heintze–Karcher inequality as a *synthetic curvature characterization* in the Finsler \(CD(K,N)\)-framework, unifying Riemannian, weighted, and anisotropic settings under a single integral inequality principle?

Why is it a "good" problem:  
This formulation resolves the earlier ambiguity by correctly imposing the **Finsler curvature lower bound**—the Ohta–Sturm \(Ric_N^F\ge (N-1)K\)—which governs volume distortion in Finsler geometry. The problem therefore cleanly ties the conjectured inequality to the intrinsic curvature of the Finsler metric, not an auxiliary Riemannian one.  

The question lies precisely at the interface of three deep theories:
1. **Comparison geometry:** It seeks a curvature–sharp integral inequality controlling hypersurface geometry by curvature bounds.  
2. **Finsler and anisotropic analysis:** It generalizes from isotropic to anisotropic curvature flows and volume measures, extending beyond known Minkowski-type results.  
3. **Synthetic geometry:** It asks whether such an inequality characterizes curvature–dimension conditions \(CD(K,N)\) in the nonsmooth, possibly non-reversible Finsler realm.

The problem is expressed in minimal but geometrically loaded terms—an inequality involving curvature, mean curvature, and volume—but embodies profound analytical and structural challenges.  
Success would connect **Heintze–Karcher principles**, **Ohta–Sturm Finsler curvature–dimension theory**, and **anisotropic geometric measure theory**, possibly yielding a unifying bridge across Riemannian, weighted, and Finsler geometries. This makes it a paradigmatic “good” research problem: conceptually simple, logically precise, cross-disciplinary, and open.